\chapter{Псевдокод алгоритмов маршрутизации}\label{appendix-algorithms}
\addcontentsline{toc}{chapter}{Приложение Б. Псевдокод алгоритмов маршрутизации}

Данное приложение содержит формальное описание ключевых алгоритмов, разработанных в рамках исследования. Для оформления псевдокода используется пакет algorithm2e.

\section{Основной алгоритм маршрутизации}

\begin{algorithm}
\SetAlgoLined
\KwData{Сеть $G=(V,E)$, абонент-источник $s$, абонент-назначения $d$}
\KwResult{Оптимальный маршрут $path$}
\BlankLine
Инициализировать таблицу маршрутов $RT \gets \emptyset$\;
Обнаружить соседей $N \gets \text{GetNeighbors}(s)$\;
\For{каждого $n_i \in N$}{
    Рассчитать метрику качества связи:\;
    $q_i \gets \alpha \cdot \text{RSSI}(n_i) + (1-\alpha)\cdot\frac{E_{res}(n_i)}{E_{init}(n_i)}$\;
    Добавить запись в $RT$: $(n_i, q_i, \text{HopCount}=1)$\;
}
\While{$RT$ не содержит маршрут к $d$}{
    Выбрать промежуточный абонент $v$ с максимальной $q_i$\;
    Обновить $RT$ через $v$:\;
    $RT \gets RT \cup \text{DiscoverRoutes}(v, d)$\;
}
Найти маршрут с минимальной метрикой:\;
$path \gets \text{маршрут } p \in RT \text{ с минимальным } \left(\sum_{v \in p} \frac{1}{q_v} + \lambda \cdot \text{HopCount}(p)\right)$\;
\caption{Адаптивный алгоритм маршрутизации}\label{alg:adaptive-routing}
\end{algorithm}

\section{Алгоритм прогнозирования позиций}

\begin{algorithm}
\SetAlgoLined
\KwData{Текущие координаты абонентов $\mathbf{x}_t$, скорости $\mathbf{v}_t$}
\KwResult{Прогнозируемые координаты $\mathbf{x}_{t+\Delta t}$}
\BlankLine
Инициализировать матрицу состояния:\;
$\mathbf{F} \gets \begin{bmatrix}1 & \Delta t \\ 0 & 1\end{bmatrix}$\;
$\mathbf{Q} \gets \begin{bmatrix}0.1 & 0 \\ 0 & 0.1\end{bmatrix}$\;
\For{каждого абонента $i$}{
    Предсказание состояния:\;
    $\mathbf{x}_i^{pred} \gets \mathbf{F} \mathbf{x}_i$\;
    $\mathbf{P}_i^{pred} \gets \mathbf{F} \mathbf{P}_i \mathbf{F}^T + \mathbf{Q}$\;
    \BlankLine
    Коррекция по измерениям:\;
    $\mathbf{K} \gets \mathbf{P}_i^{pred} (\mathbf{P}_i^{pred} + \mathbf{R})^{-1}$\;
    $\mathbf{x}_i \gets \mathbf{x}_i^{pred} + \mathbf{K}(\mathbf{z}_i - \mathbf{x}_i^{pred})$\;
    $\mathbf{P}_i \gets (\mathbf{I} - \mathbf{K})\mathbf{P}_i^{pred}$\;
}
\caption{Прогнозирование позиций абонентов}\label{alg:position-prediction}
\end{algorithm}

\section{Алгоритм управления энергией}

\begin{algorithm}
\SetAlgoLined
\KwData{Текущий заряд батареи $E_{curr}$, качество связи $q$}
\KwResult{Параметры энергосбережения}
\BlankLine
\uIf{$E_{curr} < 0.2 \cdot E_{max}$}{
    Установить режим энергосбережения:\;
    $T_{sleep} \gets T_{max}$\;
    $P_{tx} \gets P_{min}$\;
}
\Else{
    Адаптивная настройка:\;
    $T_{sleep} \gets T_{base} + k \cdot \frac{E_{curr}}{E_{max}}$\;
    $P_{tx} \gets P_{min} + (P_{max}-P_{min})\cdot\frac{q}{q_{max}}$\;
}
\caption{Адаптивное управление энергией}\label{alg:energy-management}
\end{algorithm}

\section{Примеры расчетов}

Для иллюстрации работы алгоритмов приведем пример расчета метрики маршрута:

\begin{equation}
\text{Metric} = 0.7 \cdot \frac{85}{100} + 0.3 \cdot \frac{120}{150} = 0.77
\end{equation}

где:
\begin{itemize}
\item 85 - текущий уровень сигнала (RSSI)
\item 100 - максимальный RSSI
\item 120 - остаточная энергия абонента (мДж)
\item 150 - начальная энергия абонента (мДж)
\end{itemize}